% theme: vector_magnitude_and_component_applications
\MajorQuestionLesson{ベクトルの大きさと成分の利用}
\SetRowSpace{5mm}

\TypeQuestionSingle{次のベクトルの大きさを求めなさい。}{component_magnitude}
\QQRowThree{$\left(\dfrac32,2\right)$}{}{$(-5,12)$}{}{$(1,-1)$}{}
\QQRowThree{$(-6,-8)$}{}{$(2,\sqrt5)$}{}{$(-4,7)$}{}

\TypeQuestionSingle{次の条件を満たす実数 $x$ を求めなさい。}{magnitude_equation}
\QQRow{$\left|\left(x,\dfrac32\right)\right|=\dfrac52$}{}{$\left|(3,x)\right|=\sqrt{34}$}{}
\QQRow{$\left|(x-1,2)\right|=\sqrt{13}$}{}{$\left|(2x,-3)\right|=5$}{}

\TypeQuestionSingle{次のベクトルと同じ向きの単位ベクトルを求めなさい。}{unit_vector}
\QQRowThree{$\left(\dfrac32,-2\right)$}{}{$(-8,6)$}{}{$(1,-\sqrt3)$}{}
\QQRowThree{$(-5,-12)$}{}{$(0,7)$}{}{$(2,2)$}{}

\LessonPageBreak
\SetRowSpace{5mm}

\TypeQuestionSingle{次の2つのベクトルが平行になるように、実数 $k$ の値を求めなさい。}{parallel_condition_components}
\QQRow{$\left(\dfrac12,\dfrac32\right),\ (k,6)$}{}{$(-4,6),\ (2,k)$}{}
\QQRow{$(k-1,2),\ (3,k+1)$}{}{$(k+1,2k-1),\ (3,5)$}{}

\TypeQuestionSingle{$\vec{a}=(2,1)$、$\vec{b}=(-1,2)$ とする。次のベクトルを $s\vec{a}+t\vec{b}$ の形で表しなさい。}{component_linear_combination}
\QQRowThree{$\vec{c}=\left(\dfrac32,2\right)$}{}{$\vec{d}=(-4,3)$}{}{$\vec{e}=(0,-5)$}{}

\TypeQuestionSingle{次の問いに答えなさい。}{component_applications}
\QQ{3点 $\mathrm{A}(-1,2)$、$\mathrm{B}(5,-1)$、$\mathrm{C}(2,5)$ に対して、$2\overrightarrow{\mathrm{AB}}=3\overrightarrow{\mathrm{CD}}$ を満たす点 $\mathrm{D}$ の座標を求めなさい。}{}
\QQ{$\vec{a}=(4,-2)$、$\vec{b}=(1,1)$ とする。$\vec{p}=\vec{a}+t\vec{b}$ の大きさの最小値と、そのときの実数 $t$ の値を求めなさい。}{}
\QQ{4点 $\mathrm{A}(1,-2)$、$\mathrm{B}(5,1)$、$\mathrm{C}(x,y)$、$\mathrm{D}(-2,4)$ をこの順に結んで平行四辺形ができるとき、$x,y$ の値を求めなさい。}{}
\QQ{3点 $\mathrm{A}(1,1)$、$\mathrm{B}(5,2)$、$\mathrm{C}(2,5)$ がある。これらのうち3点を頂点とする平行四辺形の第4の頂点 $\mathrm{D}$ の座標をすべて求めなさい。}{}
